\section{Spreading, forward error, and the additive schedule}\label{sec:spread-and-error}

\subsection{Filter estimates and modulo-\texorpdfstring{$r$}{r} folding}\label{sec:filter-norms}

At recursion depth $\ell$, let $R_\ell = 2^{K_\ell} = t_{1,\ell} r_\ell$, with working cutoff $p_{\mathrm{work},\ell}$ and child precision $p_{\ell+1}$. We distinguish four filter objects for each row $j_1 \in \{0, \dots, t_{1,\ell}-1\}$.
\begin{enumerate}[label=(\alph*),leftmargin=*]
  \item The \emph{unfolded sample sequence} is $\kappa^{(j_1)}_d := g(d + j_1/t_{1,\ell})$ for $d \in \Z$.
  \item The \emph{exact $r_\ell$-periodic filter} is $h^{(j_1),\mathrm{per}}_d := \sum_{a \in \Z} \kappa^{(j_1)}_{d + a r_\ell}$ for $d \in \{0, \dots, r_\ell-1\}$.
  \item The \emph{truncated modulo-$r_\ell$ folded filter} is $h^{(j_1),\mathrm{tr}}_d := \sum_{|e| \le m_{1,\ell}, \, e \equiv d \pmod{r_\ell}} \kappa^{(j_1)}_e$, supported on at most $\tilde{W}_\ell := \min(W_\ell, r_\ell)$ indices where $W_\ell := 2 m_{1,\ell} + 1$.
  \item The \emph{guarded $q_{\mathrm{child}}$-grid tap vector} $\bar{H}^{(j_1)} \in \Rg_{q_{\mathrm{child}}}^{\tilde{W}_\ell}$ and its \emph{re-indexed $r_\ell$-periodic filter} $\bar{h}^{(j_1)} \in \Rg_{q_{\mathrm{child}}}^{r_\ell}$ are formed by evaluating $\kappa^{(j_1)}_e$ ($|e| \le m_{1,\ell}$) via Lemma~\ref{lem:elementary} at guard precision $v_{\mathrm{tap}} := q_{\mathrm{child}} + \lceil \log_2(W_\ell + 1) \rceil + 10$, summing residue classes $a = (-e) \bmod r_\ell$ when $W_\ell > r_\ell$ (giving cyclic offset $o := 0$, or reversing the $W_\ell$ contiguous taps $\kappa^{(j_1)}_{m_{1,\ell} - a}$ with cyclic offset $o := m_{1,\ell}$ when $W_\ell \le r_\ell$), rounding each tap once via $\Quant_{q_{\mathrm{child}}}$ to obtain $\bar{H}^{(j_1)}_a$, and setting $\bar{h}^{(j_1)}_d := \bar{H}^{(j_1)}_a$ if $d \equiv o - a \pmod{r_\ell}$ for $a \in \{0, \dots, \tilde{W}_\ell - 1\}$ and $\bar{h}^{(j_1)}_d := 0$ otherwise.
\end{enumerate}

The proof of Lemma~\ref{lem:filter-norms} is given in Appendix~\ref{app:filter-norms-proof}.

\begin{lemma}[Filter estimates]\label{lem:filter-norms}
Let $p_{\mathrm{work},\ell} \ge 8$, $\alpha_\ell = 2\sqrt{p_{\mathrm{work},\ell} \ln 2 / \pi}$, $p_{\ell+1} \ge \max(p_{\mathrm{work},\ell}, 20)$, and $m_{1,\ell} \ge \left\lceil \alpha_\ell\sqrt{p_{\ell+1} \ln 2 / \pi} \right\rceil$. For every row $j_1 \in \{0, \dots, t_{1,\ell}-1\}$:
\begin{enumerate}[label=\textbf{\arabic*.},leftmargin=*]
  \item \textbf{Unfolded $\ell_1$ bound:}
  \begin{equation}\label{eq:unfolded-l1}
    \begin{aligned}
      \|\kappa^{(j_1)}\|_1 = \sum_{d \in \Z} \left|g\!\left(d + \frac{j_1}{t_{1,\ell}}\right)\right|
      &\le \frac{2}{\pi}\ln\alpha_\ell + \frac{5.25 - \ln\pi + E_1(1)}{\pi} \\
      &\le \frac{2}{\pi}\ln\alpha_\ell + \frac{5.25 - \ln\pi + \frac{7}{6e}}{\pi} < \frac{2}{\pi}\ln\alpha_\ell + 1.444 < \frac{2}{\pi}\ln\alpha_\ell + 1.62,
    \end{aligned}
  \end{equation}
  where $E_1(1) = \int_1^\infty e^{-t} t^{-1}\,dt < 0.2194$.
  \item \textbf{Spatial tail bound:}
  \begin{equation}\label{eq:spatial-tail}
    \tau := \sum_{|d| > m_{1,\ell}} \left|g\!\left(d + \frac{j_1}{t_{1,\ell}}\right)\right| \le \left(\frac{1}{\pi m_{1,\ell}} + \frac{1}{\pi p_{\ell+1} \ln 2}\right) 2^{-p_{\ell+1}} < 0.077 \cdot 2^{-p_{\ell+1}} < 0.10 \cdot 2^{-p_{\ell+1}}.
  \end{equation}
  \item \textbf{Periodic and folded filter bounds:}
  \begin{equation}\label{eq:periodic-norms}
    \begin{aligned}
      \|h^{(j_1),\mathrm{per}}\|_2 &\le 1, \qquad & \|h^{(j_1),\mathrm{per}}\|_1 &\le \sqrt{r_\ell}, \\
      \|h^{(j_1),\mathrm{tr}} - h^{(j_1),\mathrm{per}}\|_1 &\le \tau, \qquad & \|\bar{h}^{(j_1)} - h^{(j_1),\mathrm{tr}}\|_1 &\le \tilde{W}_\ell 2^{-p_{\ell+1}-1}.
    \end{aligned}
  \end{equation}
\end{enumerate}
\end{lemma}

\subsection{Spreading by overlap-add and Kronecker substitution}\label{sec:kronecker}

To evaluate each row's length-$r_\ell$ cyclic convolution with its short filter $\bar{H}^{(j_1)}$ of length $\tilde{W}_\ell \le r_\ell$, we combine classical overlap-add block convolution~\cite{Nussbaumer82} with Kronecker substitution for polynomial multiplication~\cite{Schonhage82,Harvey09} in a signed two's-complement radix representation.

\begin{lemma}[Signed radix Kronecker packing and extraction]\label{lem:kronecker}
Let $s \ge 1$ and $B_{\mathrm{slot}} \ge 4$ be integers, and let $a_0, \dots, a_{s-1} \in \Z$ satisfy $|a_j| < 2^{B_{\mathrm{slot}}-1}$. Define the signed radix-$2^{B_{\mathrm{slot}}}$ evaluation
\begin{equation}\label{eq:pack-def}
  \Pack_{B_{\mathrm{slot}}}(a_0, \dots, a_{s-1}) := \sum_{j=0}^{s-1} a_j \, 2^{j B_{\mathrm{slot}}} \in \Z.
\end{equation}
\begin{enumerate}[label=(\roman*),leftmargin=*]
  \item Given the signed binary representations of $a_0, \dots, a_{s-1}$, the integer $\Pack_{B_{\mathrm{slot}}}(a)$ has bit-length at most $s B_{\mathrm{slot}} + 1$ and can be assembled in $O(s B_{\mathrm{slot}})$ bit operations on a multitape Turing machine via a single low-to-high streaming carry pass.
  \item Conversely, if $c_0, \dots, c_{m-1} \in \Z$ satisfy $|c_k| < 2^{B_{\mathrm{slot}}-1}$ for all $k$ and $T_0 := \sum_{k=0}^{m-1} c_k 2^{k B_{\mathrm{slot}}}$, then the recurrence
  \[
    r_k := T_k \bmod 2^{B_{\mathrm{slot}}} \in [0, 2^{B_{\mathrm{slot}}}), \quad c_k := \begin{cases} r_k, & r_k < 2^{B_{\mathrm{slot}}-1}, \\ r_k - 2^{B_{\mathrm{slot}}}, & r_k \ge 2^{B_{\mathrm{slot}}-1}, \end{cases} \quad T_{k+1} := \frac{T_k - c_k}{2^{B_{\mathrm{slot}}}},
  \]
  uniquely recovers $c_0, \dots, c_{m-1}$ in $O(m B_{\mathrm{slot}})$ bit operations.
\end{enumerate}
\end{lemma}

\begin{proof}
For (i), write each $a_j \in (-2^{B_{\mathrm{slot}}-1}, 2^{B_{\mathrm{slot}}-1})$ with a borrow $\sigma_j \in \{-1, 0\}$ into the next slot: initialize $\sigma_0 := 0$, and for $j = 0, \dots, s-1$, let $a'_j := a_j + \sigma_j$. If $a'_j \ge 0$, emit the $B_{\mathrm{slot}}$-bit unsigned word $u_j := a'_j \in [0, 2^{B_{\mathrm{slot}}-1})$ and set $\sigma_{j+1} := 0$; if $a'_j < 0$, emit the $B_{\mathrm{slot}}$-bit unsigned word $u_j := a'_j + 2^{B_{\mathrm{slot}}} \in [2^{B_{\mathrm{slot}}-1}, 2^{B_{\mathrm{slot}}})$ and set $\sigma_{j+1} := -1$. Then $\sum_{j=0}^{s-1} a_j 2^{j B_{\mathrm{slot}}} = \sigma_s 2^{s B_{\mathrm{slot}}} + \sum_{j=0}^{s-1} u_j 2^{j B_{\mathrm{slot}}}$, where concatenating the $B_{\mathrm{slot}}$-bit words $u_j$ takes $O(s B_{\mathrm{slot}})$ bit operations. For (ii), since $T_k = c_k + 2^{B_{\mathrm{slot}}} T_{k+1}$, we have $c_k \equiv T_k \pmod{2^{B_{\mathrm{slot}}}}$, and because $[-2^{B_{\mathrm{slot}}-1}, 2^{B_{\mathrm{slot}}-1})$ contains a unique representative of each residue class modulo $2^{B_{\mathrm{slot}}}$, $c_k$ and $T_{k+1}$ are uniquely determined and extracted in $O(B_{\mathrm{slot}})$ bit operations per slot.
\end{proof}

\begin{theorem}[Spreading by overlap-add and Kronecker substitution]\label{thm:spread}
At any non-leaf node ($R_\ell > 4 p_\ell^2$, with $p_\ell \ge 8$) under the hypotheses of Lemma~\ref{lem:filter-norms} with $p_{\mathrm{work},\ell} = p_\ell + K_\ell + 6$ and $p_{\ell+1} = p_\ell + 2 K_\ell + 12$, let $x \in \Cg_q^{R_\ell}$ satisfy $\|x\|_\infty \le 2^{p_\ell}$, let $q := q_{\mathrm{child}} = p_{\ell+1} + 2$, $\tilde{W}_\ell := \min(2m_{1,\ell}+1, r_\ell) \le r_\ell$, and
\begin{equation}\label{eq:slot-width}
  B_{\mathrm{slot}} := 2q + p_\ell + \lceil \log_2(\tilde{W}_\ell + 1) \rceil + 4 = \Theta(p_{\ell+1}).
\end{equation}
Computing the cyclic row convolutions for all $j_1 \in \{0, \dots, t_{1,\ell}-1\}$ with the guarded $q$-grid filters $\bar{H}^{(j_1)}$ via overlap-add of block length $\tilde{W}_\ell$, signed Kronecker substitution at slot width $B_{\mathrm{slot}}$, exact integer overlap addition on $\Cg_{2q}$, and a single final rounding per output entry to $\Cg_q$ produces an array $\widehat{U} \in \Cg_q^{t_{1,\ell} \times r_\ell}$ satisfying
\begin{equation}\label{eq:spread-err}
  \|\widehat{U} - U^\circ\|_\infty \le \bigl(\tfrac{1}{2}\tilde{W}_\ell + 0.10\bigr) 2^{-p_{\ell+1}}\|x\|_\infty + \frac{\sqrt{2}}{8} 2^{-p_{\ell+1}} \le 2 r_\ell \, 2^{-p_{\ell+1}} \max(1, \|x\|_\infty)
\end{equation}
in time
\begin{equation}\label{eq:spread-cost}
  \mathcal{C}_{\mathrm{spread}}(R_\ell, p_{\ell+1}) = O\!\left(\frac{R_\ell}{\tilde{W}_\ell} \, \mathsf{M}\bigl(\tilde{W}_\ell B_{\mathrm{slot}} + 1\bigr) + R_\ell p_{\ell+1} \log_2 t_{1,\ell}\right)
\end{equation}
bit operations, where $\mathsf{M}(N)$ denotes the bit complexity of multiplying two $N$-bit integers and $\tilde{W}_\ell B_{\mathrm{slot}} + 1 = \Theta(\tilde{W}_\ell p_{\ell+1})$ is the operand length of the block products in the proof.
\end{theorem}

\begin{proof}
First, by Lemma~\ref{lem:routing}(i), de-interleaving $x \in \Cg_q^{R_\ell}$ into the $t_{1,\ell}$ contiguous rows $v^{(j_1)} = \bigl(x(j_1 + t_{1,\ell} j_2)\bigr)_{j_2=0}^{r_\ell-1}$ takes $O(R_\ell p_{\ell+1} \log_2 t_{1,\ell})$ bit operations.
At any non-leaf node, $R_\ell > 4 p_\ell^2 \ge 256$ (since $p_\ell \ge 8$), so $K_\ell \ge 9$, $r_\ell = 2^{\lceil K_\ell / 2 \rceil} \ge 32$, and $p_\ell < \frac{1}{2}\sqrt{R_\ell} \le \frac{1}{2} r_\ell$. Under $p_{\mathrm{work},\ell} = p_\ell + K_\ell + 6$, since $K_\ell \le 2\log_2 r_\ell < r_\ell$ and $6 < \frac{1}{2} r_\ell$ for all $r_\ell \ge 32$, we have $p_{\mathrm{work},\ell} < 2 r_\ell$, and therefore
\begin{align*}
  \alpha_\ell &= 2\sqrt{p_{\mathrm{work},\ell}\ln 2 / \pi} < 2\sqrt{r_\ell} < r_\ell, \\
  r - \frac{2}{\pi}\ln r &\ge 32 - \frac{2}{\pi}\ln 32 > 29 > 1.62 \quad (r \ge 32), \\
  \frac{2}{\pi}\ln\alpha_\ell + 2.62 &< \frac{2}{\pi}\ln r_\ell + 2.62 \le r_\ell + 1.
\end{align*}
Meanwhile, since $p_{\ell+1} \ge p_{\mathrm{work},\ell}$ and $\sqrt{p_{\mathrm{work},\ell}\ln 2 / \pi} = \frac{1}{2}\alpha_\ell$, Lemma~\ref{lem:filter-norms} gives $m_{1,\ell} \ge \max\bigl(6, \lceil \frac{1}{2}\alpha_\ell^2 \rceil\bigr)$, so $2m_{1,\ell} + 2 \ge \max(14, \alpha_\ell^2 + 2) > \frac{2}{\pi}\ln\alpha_\ell + 2.62$ for all $\alpha_\ell > 0$. Combining both branches yields $\frac{2}{\pi}\ln\alpha_\ell + 2.62 \le \min(2m_{1,\ell}+1, r_\ell) + 1 = \tilde{W}_\ell + 1$ across all non-leaves.

For each row $j_1$, partition $v^{(j_1)} \in \Cg_q^{r_\ell}$ into $\lceil r_\ell / \tilde{W}_\ell \rceil \le \frac{2 r_\ell}{\tilde{W}_\ell}$ contiguous blocks of length $L_b \le \tilde{W}_\ell$. On $\Cg_q$, each scaled input coordinate $2^q \operatorname{Re}(v^{(j_1)}_{j_2}), 2^q \operatorname{Im}(v^{(j_1)}_{j_2}) \in \Z$ has absolute value $\le 2^{q + p_\ell}$, and by Lemma~\ref{lem:filter-norms}(1),(3), the guarded filter satisfies
\[
  \|\bar{H}^{(j_1)}\|_1 = \|\bar{h}^{(j_1)}\|_1 \le \|h^{(j_1),\mathrm{tr}}\|_1 + \frac{1}{2}\tilde{W}_\ell 2^{-p_{\ell+1}} \le \|\kappa^{(j_1)}\|_1 + 1 < \frac{2}{\pi}\ln\alpha_\ell + 2.62 \le \tilde{W}_\ell + 1,
\]
so $\sum_{a=0}^{\tilde{W}_\ell-1} |2^q \bar{H}^{(j_1)}_a| \le (\tilde{W}_\ell + 1) 2^q$. By the $\ell_\infty$--$\ell_1$ inequality, every integer linear-convolution coefficient between a block and $2^q \bar{H}^{(j_1)}$ has absolute value at most
\[
  2^{q+p_\ell} \cdot (\tilde{W}_\ell + 1) 2^q = (\tilde{W}_\ell + 1) 2^{2q + p_\ell} \le 2^{B_{\mathrm{slot}}-4} < 2^{B_{\mathrm{slot}}-1}
\]
unconditionally (whether $W_\ell \le r_\ell$ or $W_\ell > r_\ell$).

By Lemma~\ref{lem:kronecker}, packing the real and imaginary parts of each block into $\Pack_{B_{\mathrm{slot}}}$, multiplying each by the pre-packed filter integer $\Pack_{B_{\mathrm{slot}}}(2^q \bar{H}^{(j_1)})$ (replayed from $T_{\mathrm{tab}}$ for each block, performing two integer multiplications of bit-length $\le \tilde{W}_\ell B_{\mathrm{slot}} + 1 = \Theta(\tilde{W}_\ell p_{\ell+1})$ per block), and extracting the $L_b + \tilde{W}_\ell - 1$ coefficients recovers the exact integer block convolutions at fractional precision $2q$. To accumulate the overlapping block outputs modulo $r_\ell$ on a multitape Turing machine, as each block $b$ of row $j_1$ is unpacked into $L_b + \tilde{W}_\ell - 1 \ge \tilde{W}_\ell$ coefficients, the $(\tilde{W}_\ell - 1)$-entry carry-over tail from block $b - 1$ (held on scratch tape $T_1$, initialized to $0$ at $b = 0$) is added in lockstep to its first $\tilde{W}_\ell - 1$ coefficients. The first $L_b$ resulting coefficients are appended sequentially to scratch tape $T_2$, while the remaining $\tilde{W}_\ell - 1$ coefficients overwrite $T_1$. After the final block of row $j_1$, the residual $(\tilde{W}_\ell - 1)$-entry tail on $T_1$ wraps cyclically modulo $r_\ell$ to row positions $0, \dots, \tilde{W}_\ell - 2$: rewinding $T_2$ to the start of row $j_1$ and $T_1$ to index $0$, one pass over the first $\tilde{W}_\ell - 1$ entries adds this wrapped tail in $O(\tilde{W}_\ell B_{\mathrm{slot}}) = O(r_\ell p_{\ell+1})$ steps.

Finally, when $W_\ell \le r_\ell$, the reversed contiguous filter $\kappa^{(j_1)}_{m_{1,\ell} - a}$ ($0 \le a < W_\ell$) offsets the cyclic convolution indices by $o = m_{1,\ell} \in \{0, \dots, r_\ell - 1\}$ (whereas when $W_\ell > r_\ell$, the modulo-$r_\ell$ folded filter is already indexed by $a = (-e) \bmod r_\ell$ with offset $o = 0$). During the single final rounding pass from scale $2q$ to $\Cg_q$, copying the first $o$ entries of $T_2$ to $T_1$ and streaming the remaining $r_\ell - o$ entries of $T_2$ followed by the $o$ entries of $T_1$ rotates row $j_1$ into natural index order $k_2 = 0, \dots, r_\ell - 1$ in $O(r_\ell w_\ell) = O(r_\ell p_{\ell+1})$ steps while incurring a single rounding error $\le \frac{\sqrt{2}}{2} 2^{-q} = \frac{\sqrt{2}}{8} 2^{-p_{\ell+1}}$ per entry. Combining this with Lemma~\ref{lem:filter-norms}(3) ($\|\bar{h}^{(j_1)} - h^{(j_1),\mathrm{per}}\|_1 \le (\frac{1}{2}\tilde{W}_\ell + 0.10) 2^{-p_{\ell+1}}$) yields \eqref{eq:spread-err}. Across all $t_{1,\ell}$ rows, there are at most $2 R_\ell / \tilde{W}_\ell$ blocks, proving \eqref{eq:spread-cost}.
\end{proof}

\begin{remark}[Modulo-$r_\ell$ folding at small non-leaves]\label{rem:folding}
At any non-leaf $R_\ell > 4 p_\ell^2$, we have $p_\ell < \frac{1}{2}\sqrt{R_\ell} \le \frac{1}{2} r_\ell$, so $W_\ell = 2m_{1,\ell} + 1 = \Theta(p_{\ell+1}) = O(r_\ell)$. For the default radius $m_{1,\ell} = \lceil \alpha_\ell\sqrt{p_{\ell+1}\ln 2 / \pi} \rceil$ under schedule~\eqref{eq:cert-sched}, $W_\ell < r_\ell$ at every non-leaf reached from the root parameters (for $n > 64$ one has $p_\ell \ge 2K_\ell + 14$ at every depth, whence $W_\ell \le 1.53\,p_\ell + 3 < 2p_\ell < r_\ell$), but $W_\ell > r_\ell$ does occur inside the hypotheses of Theorem~\ref{thm:spread} (for the larger admissible radii, up to $m_{1,\ell} = p_{\ell+1}$, and for parameter pairs not reached from any root) and outside them (for the lower leaf cutoff of the C~program, Section~\ref{sec:empirical}): for example, the parameters $(p_\ell, K_\ell) = (15, 10)$, which satisfy $R_\ell = 1024 > 4 p_\ell^2 = 900$ but are not reached from any root, give with $m_{1,\ell} = \lceil \alpha_\ell\sqrt{p_{\ell+1}\ln 2 / \pi} \rceil = 17$ the values $r_\ell = 32$ and $W_\ell = 35 > r_\ell$. Folding modulo $r_\ell$ (and never modulo $W_\ell$) preserves the exact cyclic convolution modulo $r_\ell$ while keeping the effective filter length $\tilde{W}_\ell = \min(W_\ell, r_\ell) \le r_\ell$.
\end{remark}

\subsection{Forward error and the additive schedule}\label{sec:forward-error}

\begin{lemma}[Forward error and the additive schedule]\label{lem:forward}
At recursion depth $\ell$, let $R_\ell = 2^{K_\ell}$ ($K_\ell \ge 2$) and let the target accuracy be $p_\ell \ge \max(8, K_\ell)$ on an input $x \in \Cg_{p_\ell+2}^{R_\ell}$ with $\|x\|_\infty \le 2^{p_\ell}$. Define the additive precision schedule
\begin{equation}\label{eq:cert-sched}
  p_{\mathrm{work},\ell} := p_\ell + K_\ell + 6, \qquad p_{\ell+1} := p_{\mathrm{work},\ell} + K_\ell + 6 = p_\ell + 2 K_\ell + 12, \qquad q_{\mathrm{child}} := p_{\ell+1} + 2,
\end{equation}
and choose any integer window radius $m_{1,\ell} \in \bigl[\lceil \alpha_\ell\sqrt{p_{\ell+1}\ln 2 / \pi} \rceil, \, p_{\ell+1}\bigr]$ (for instance $m_{1,\ell} := \lceil \alpha_\ell\sqrt{p_{\ell+1}\ln 2 / \pi} \rceil$ or $m_{1,\ell} := p_{\ell+1}$). Assume inductively that every recursive child call at length $r_\ell$ and precision $p_{\ell+1}$ satisfies $\|\widehat{C}_{v_1} - F_{r_\ell}(\widehat{V}_{v_1, \cdot})\|_\infty \le 2^{-p_{\ell+1}}\max(1, \|\widehat{V}_{v_1, \cdot}\|_\infty)$. Then:
\begin{enumerate}[label=\textbf{\arabic*.},leftmargin=*]
  \item \textbf{Forward error progression at a non-leaf:} Assume $R_\ell > 4 p_\ell^2$ (non-leaf), and let $\mathcal{A} := \max(1, \|x\|_\infty)$ and $u := 2^{-p_{\ell+1}}$ (following the three-step decomposition of \texttt{FastDFT} in Algorithm~\ref{alg:main} and Listing~\ref{lst:multiply}: Step~1 row spreading, Step~2 polynomial-transform 2D DFT via Steps~2a--2d, and Step~3 column deconvolution). Let $Z := F_{R_\ell} x$ be the exact 1D DFT, $Z^*$ be the exact solution of the retained two-term system with ideal right-hand side $Y^\circ = F_{t_{1,\ell}, r_\ell} U^\circ$, $Z^\dagger$ be the exact solution of the retained two-term system with the routed child output $\widehat{Y}$, and $Z^{\mathrm{comp}}$ be the output of \texttt{CyclicSolve}/\texttt{NyquistSolve}. Then:
  \begin{equation}\label{eq:forward-split}
    \begin{aligned}
      \|Z^{\mathrm{comp}} - F_{R_\ell} x\|_\infty
      &\le \underbrace{\frac{1}{2} t_{1,\ell} \, 2^{-p_{\mathrm{work},\ell}} \mathcal{A}}_{\|Z^* - Z\|_\infty \text{ (aliases)}} + \underbrace{27 R_\ell \, u \mathcal{A}}_{\|Z^\dagger - Z^*\|_\infty \text{ (Steps 1--2)}} + \underbrace{160 R_\ell \, u \mathcal{A}}_{\|Z^{\mathrm{comp}} - Z^\dagger\|_\infty \text{ (Step 3)}} \\
      &\le \frac{1}{2} t_{1,\ell} \, 2^{-p_{\mathrm{work},\ell}} \mathcal{A} + 200 R_\ell \, u \mathcal{A}.
    \end{aligned}
  \end{equation}
  \item \textbf{Inductive closure at non-leaves and leaves:} After quantizing $Z^{\mathrm{comp}}$ to $\Cg_{p_\ell + 2}$ and transposing to natural index order when $R_\ell > 4 p_\ell^2$ (or after executing \texttt{BaseRadix2FFT} at internal precision $q_b := p_\ell + K_\ell + 4$, which \texttt{FastDFT} does when $R_\ell \le 4 p_\ell^2$; the \texttt{BaseRadix2FFT} bound below holds for every length $R_\ell = 2^{K_\ell}$, $K_\ell \ge 2$, and is also used at the root length by the bottom stage of Section~\ref{sec:cost-proof}), the returned vector $\widehat{Z} = \texttt{FastDFT}(x, R_\ell, p_\ell)$ satisfies
  \begin{equation}\label{eq:inductive-bound}
    \|\widehat{Z} - F_{R_\ell} x\|_\infty \le 2^{-p_\ell} \max(1, \|x\|_\infty).
  \end{equation}
\end{enumerate}
\end{lemma}

\begin{proof}
First verify that the interval $\bigl[\lceil \alpha_\ell\sqrt{p_{\ell+1}\ln 2 / \pi} \rceil, \, p_{\ell+1}\bigr]$ is non-empty and satisfies the hypotheses $p_{\mathrm{work},\ell} \ge 16 \ge 8$ and $p_{\ell+1} \ge 24 \ge \max(p_{\mathrm{work},\ell}, 20)$ of Lemma~\ref{lem:filter-norms}. Since $\alpha_\ell = 2\sqrt{p_{\mathrm{work},\ell}\ln 2 / \pi}$ and $\frac{2\ln 2}{\pi} < 0.442 < 1$, for all $p_\ell \ge 8$ and $K_\ell \ge 2$ we have
\[
  \alpha_\ell\sqrt{p_{\ell+1}\ln 2 / \pi} = \frac{2\ln 2}{\pi}\sqrt{p_{\mathrm{work},\ell}\,p_{\ell+1}} < 0.442\,p_{\ell+1} < p_{\ell+1}.
\]

\textbf{Steps 1 and 2 of Algorithm~\ref{alg:main} (with Step~2 realized by the five steps of Proposition~\ref{thm:twisted}): $Y^\circ \to \widehat{Y}$ at a non-leaf $R_\ell > 4 p_\ell^2$.}
By Lemma~\ref{lem:filter-norms}(3), $\|U^\circ\|_\infty \le \max_{j_1} \|h^{(j_1),\mathrm{per}}\|_1 \|x\|_\infty \le \sqrt{r_\ell}\mathcal{A}$, and by Theorem~\ref{thm:spread} (applicable since $R_\ell > 4 p_\ell^2$), $\|\widehat{U} - U^\circ\|_\infty \le 2 r_\ell u \mathcal{A}$.
In Step~2a (pre-twist), by \eqref{eq:faithful-table} the 1-ulp faithful table roots $\bar{\zeta}_{k_2} \in \Cg_{q_{\mathrm{child}}}$ satisfy $|\bar{\zeta}_{k_2} - \zeta^{k_2}| \le \frac{\sqrt{2}}{4} u < 0.354 u$ and $|\bar{\zeta}_{k_2}| \le 1 + 0.354 u$. Since $r_\ell \ge 2$ implies $\frac{0.354}{\sqrt{r_\ell}} + \frac{0.177}{r_\ell} < 0.34$ and $\frac{\log_2 t_{1,\ell}}{r_\ell} \le \frac{1}{2}$, each entry $\widehat{P}_{j_1, k_2} = \RM_{q_{\mathrm{child}}}(\bar{\zeta}_{k_2}, \widehat{U}_{j_1, k_2})$ differs from $P^\circ_{j_1, k_2} = \zeta^{k_2} U^\circ_{j_1, k_2}$ by at most
\[
  |\bar{\zeta}_{k_2}| \, |\widehat{U}_{j_1, k_2} - U^\circ_{j_1, k_2}| + |\bar{\zeta}_{k_2} - \zeta^{k_2}| \, |U^\circ_{j_1, k_2}| + \frac{\sqrt{2}}{8} u \le (1 + u)(2 r_\ell u \mathcal{A}) + 0.354 \sqrt{r_\ell} u \mathcal{A} + 0.177 u \mathcal{A} \le 2.6 r_\ell u \mathcal{A},
\]
and $\|\widehat{P}\|_\infty \le \sqrt{r_\ell}\mathcal{A} + 2.6 r_\ell u \mathcal{A} \le 2 r_\ell \mathcal{A}$.
In Step~2b (polynomial transform $G_{t_{1,\ell}}$), since $G_{t_{1,\ell}}$ is an average over $t_{1,\ell}$ monomial shifts and has coefficient-wise $\ell_\infty$ norm $\le 1$, Lemma~\ref{lem:syn-fft} gives $\|\widehat{Q} - Q^\circ\|_\infty \le \|\widehat{P} - P^\circ\|_\infty + 0.177 (\log_2 t_{1,\ell}) u \le 3 r_\ell u \mathcal{A}$ and $\|Q^\circ\|_\infty \le \sqrt{r_\ell}\mathcal{A}$.
In Step~2c (post-twist by $\overline{\bar{\zeta}_{k_2}}$), the same product bound and the norm chain give
\begin{align*}
  \|\widehat{V} - V^\circ\|_\infty &\le (1 + u)(3 r_\ell u \mathcal{A}) + 0.354 \sqrt{r_\ell} u \mathcal{A} + 0.177 u \mathcal{A} \le 3.6 r_\ell u \mathcal{A} < 7 r_\ell u \mathcal{A}, \\
  \|V^\circ\|_\infty &= \|Q^\circ\|_\infty \le \|P^\circ\|_\infty = \|U^\circ\|_\infty \le \sqrt{r_\ell}\mathcal{A} \le r_\ell \mathcal{A}, \\
  \|\widehat{V}_{v_1, \cdot}\|_\infty &\le r_\ell \mathcal{A} + 7 r_\ell u \mathcal{A} \le 2 r_\ell \mathcal{A} \le 2^{p_\ell + K_\ell + 1} < 2^{p_\ell + 2 K_\ell + 12} = 2^{p_{\ell+1}}.
\end{align*}
Thus each child call $\widehat{C}_{v_1} = \texttt{FastDFT}(\widehat{V}_{v_1, \cdot}, r_\ell, p_{\ell+1})$ in Step~2d satisfies the child input precondition $\|\widehat{V}_{v_1, \cdot}\|_\infty \le 2^{p_{\ell+1}}$, and by the induction hypothesis and $\|F_{r_\ell}\|_{\infty \to \infty} \le 1$,
\[
  \|\widehat{C}_{v_1} - F_{r_\ell} V^\circ_{v_1, \cdot}\|_\infty \le \|\widehat{C}_{v_1} - F_{r_\ell} \widehat{V}_{v_1, \cdot}\|_\infty + \|\widehat{V}_{v_1, \cdot} - V^\circ_{v_1, \cdot}\|_\infty \le u(2 r_\ell \mathcal{A}) + 7 r_\ell u \mathcal{A} = 9 r_\ell u \mathcal{A}.
\]
Because the odd-stride permutation \eqref{eq:odd-perm} is an index bijection, the routed array $\widehat{Y}$ satisfies $\|\widehat{Y} - Y^\circ\|_\infty \le 9 r_\ell u \mathcal{A}$, and
\begin{align*}
  \|Y^\circ\|_\infty &\le \sum_{m \in \Z} \hat{g}(u_2/r_\ell - m) \|F_{R_\ell} x\|_\infty \le \|x\|_\infty \le \mathcal{A}, \\
  \mathcal{M}_Y &= \max(1, \|\widehat{Y}\|_\infty) \le \mathcal{A} + 9 r_\ell u \mathcal{A} \le 2\mathcal{A}.
\end{align*}

\textbf{Step 3 error ($Z \to Z^* \to Z^\dagger \to Z^{\mathrm{comp}}$).}
By Theorem~\ref{thm:poisson} and Theorem~\ref{thm:deconv}(1), for every column $u_2$, the exact DFT $Z$ satisfies the retained two-term equations with right-hand side $Y^\circ - \Delta^{\mathrm{alias}}$, where $\|\Delta^{\mathrm{alias}}\|_\infty \le \Psi(0)\|F_{R_\ell} x\|_\infty < 0.15 \cdot 2^{-p_{\mathrm{work},\ell}}\mathcal{A}$ (and at $u_2 = r_\ell/2$ the alternating sum $A_{\mathrm{N}}$ in \eqref{eq:nyquist-id} is exact on $\Cg_{q_{\mathrm{child}}}$ because $q_{\mathrm{child}} = p_\ell + 2 K_\ell + 14 > p_\ell + 2 + \log_2 r_\ell$). Applying the inverse norm bound $< 2.36 \, t_{1,\ell} < 3 t_{1,\ell}$ from Theorem~\ref{thm:deconv}(2)--(3) gives
\[
  \|Z^* - Z\|_\infty \le 2.36 \, t_{1,\ell} \cdot 0.15 \cdot 2^{-p_{\mathrm{work},\ell}} \mathcal{A} < \frac{1}{2} t_{1,\ell} \, 2^{-p_{\mathrm{work},\ell}} \mathcal{A},
\]
and perturbing the right-hand side from $Y^\circ$ to $\widehat{Y}$ gives
\[
  \|Z^\dagger - Z^*\|_\infty \le 3 t_{1,\ell} \|\widehat{Y} - Y^\circ\|_\infty \le 3 t_{1,\ell} (9 r_\ell u \mathcal{A}) = 27 R_\ell u \mathcal{A}.
\]
Furthermore,
\[
  p_{\ell+1} - \lfloor K_\ell / 2 \rfloor \ge 8 + \frac{3}{2} K_\ell + 12 \ge 23 > 20, \qquad t_{1,\ell} u \le 2^{-23} < 2^{-20},
\]
so Lemma~\ref{lem:solver-err} and $\mathcal{M}_Y \le 2\mathcal{A}$ yield $\|Z^{\mathrm{comp}} - Z^\dagger\|_\infty \le 80 t_{1,\ell}^2 u (2\mathcal{A}) \le 160 R_\ell u \mathcal{A}$. Summing the three bounds proves \eqref{eq:forward-split}.

\textbf{Inductive closure.}
Under \eqref{eq:cert-sched}, $p_{\mathrm{work},\ell} - p_\ell = K_\ell + 6$ and $p_{\ell+1} - p_\ell = 2 K_\ell + 12$, with $t_{1,\ell} = 2^{\lfloor K_\ell / 2 \rfloor} \le 2^{K_\ell}$ and $R_\ell = 2^{K_\ell}$ ($K_\ell \ge 2$). Thus
\[
  \begin{aligned}
    \frac{1}{2} t_{1,\ell} \, 2^{-(p_{\mathrm{work},\ell} - p_\ell)} &\le \frac{1}{2} 2^{K_\ell} \cdot 2^{-K_\ell - 6} = 2^{-7}, \\
    200 R_\ell \, 2^{-(p_{\ell+1} - p_\ell)} &= 200 \cdot 2^{-K_\ell - 12} \le 200 \cdot 2^{-14} = \frac{25}{2048} < 2^{-6}.
  \end{aligned}
\]
Final nearest-grid rounding $\widehat{Z} = \Quant_{p_\ell + 2}(Z^{\mathrm{comp}})$ adds at most $\frac{\sqrt{2}}{8} 2^{-p_\ell} < 0.177 \cdot 2^{-p_\ell}$, so
\[
  \|\widehat{Z} - F_{R_\ell} x\|_\infty \le \left(\frac{\sqrt{2}}{8} + 2^{-7} + \frac{25}{2048}\right) 2^{-p_\ell} \mathcal{A} < 0.197 \cdot 2^{-p_\ell} \mathcal{A} < \frac{1}{4} 2^{-p_\ell} \mathcal{A} < 2^{-p_\ell} \mathcal{A}.
\]
For \texttt{BaseRadix2FFT} at any length $R_\ell = 2^{K_\ell}$, $K_\ell \ge 2$ (in particular at a leaf $R_\ell \le 4 p_\ell^2$), \texttt{BaseRadix2FFT} bit-reverses its $R_\ell = 2^{K_\ell}$ input records and executes $K_\ell = \log_2 R_\ell$ decimation-in-time radix-2 Cooley--Tukey butterfly stages on $\Cg_{q_b}$ with $q_b \ge p_\ell + K_\ell + 4$. In stage $s \in \{1, \dots, K_\ell\}$, each block of $2^s$ records is split into its upper and lower halves $(A_j)_{j=0}^{2^{s-1}-1}$ and $(B_j)_{j=0}^{2^{s-1}-1}$ on $T_1, T_2$ and paired in lockstep with the stage-$s$ 1-ulp faithful twiddle list $(\bar{\omega}_{s, j})_{j=0}^{2^{s-1}-1}$ replayed from $T_{\mathrm{tab}}$ ($|\bar{\omega}_{s, j} - \omega_{s, j}| \le \sqrt{2} \cdot 2^{-q_b}$). Each pair is then updated via $W_j := \RM_{q_b}(\bar{\omega}_{s, j}, B_j)$ and $A_j^{\mathrm{new}}, B_j^{\mathrm{new}} := \Quant_{q_b}\bigl(\frac{1}{2}(A_j \pm W_j)\bigr)$. Inductively $\|A\|_\infty, \|B\|_\infty \le 1.01\mathcal{A}$, and each stage incurs twiddle approximation error $\le \sqrt{2} \cdot 2^{-q_b}(1.01\mathcal{A})$ plus product and halving/rounding errors $\le \sqrt{2} \cdot 2^{-q_b}\mathcal{A}$, so $e_s \le (1 + 2^{-q_b})e_{s-1} + 2.85 \cdot 2^{-q_b}\mathcal{A}$. Since $K_\ell 2^{-K_\ell} \le \frac{1}{2}$ for all $K_\ell \ge 1$, across $K_\ell$ stages we obtain
\[
  e_{K_\ell} \le 2.86 K_\ell 2^{-q_b} \mathcal{A} \le 2.86 K_\ell 2^{-K_\ell - 4} 2^{-p_\ell}\mathcal{A} \le 1.43 \cdot 2^{-4} 2^{-p_\ell}\mathcal{A} < 0.090 \cdot 2^{-p_\ell}\mathcal{A}
\]
before final rounding to $\Cg_{p_\ell+2}$ ($\le 0.177 \cdot 2^{-p_\ell}\mathcal{A}$), giving $< 0.267 \cdot 2^{-p_\ell}\mathcal{A} < 2^{-p_\ell}\mathcal{A}$.
\end{proof}

\begin{remark}[Shared precision schedule]\label{rem:schedules}
Because the additive increments $p_{\mathrm{work},\ell} - p_\ell = K_\ell + 6$ and $p_{\ell+1} - p_\ell = 2 K_\ell + 12$ in \eqref{eq:cert-sched} bound the worst-case $\ell_\infty$ alias amplification ($\le 2^{-7} \cdot 2^{-p_\ell}\mathcal{A}$) and the worst-case $\ell_\infty$ rounding and child-DFT amplification ($\le 200 \cdot 2^{-K_\ell - 12} \cdot 2^{-p_\ell}\mathcal{A} < 2^{-6} \cdot 2^{-p_\ell}\mathcal{A}$) for all $K_\ell \ge 2$ with more than a factor-of-five total margin ($< 0.197 \cdot 2^{-p_\ell}\mathcal{A}$), the exact same precision schedule \eqref{eq:cert-sched} and root precision $p_0 = 2b + 2K_0 + 6$ (with window radius \mbox{$m_{1,\ell} \in [\lceil \alpha_\ell\sqrt{p_{\ell+1}\ln 2 / \pi} \rceil,\, p_{\ell+1}]$}) are used across our written proof, our Lean~4 formalization~\cite{LeanMS26} (Section~\ref{sec:lean}), and our C/GMP implementation~\cite{CGMP26} (Section~\ref{sec:empirical}). As noted in Section~\ref{sec:empirical}, the leaf threshold $R_\ell \le 4 p_\ell^2$ of the written proof and of~\cite{LeanMS26} (\texttt{isBaseLeaf}, Table~\ref{tab:lean-map}) is lowered by default to $\texttt{bt} = p_0^2$ in~\cite{CGMP26}, and, as noted in Section~\ref{sec:lean}, the tie-breaking rule on half-integer midpoints differs slightly across implementations.
\end{remark}
